\chapter{How Forge is implemented}
\label{ch:implementation}
\index{implementation}

The implementation is an OCaml library (\texttt{src/}) shared by all the
executables. Figure~\ref{fig:pipeline} shows the path of a program through
it: a common front end, then either the interpreter or the native back end.

\begin{figure}[ht]
\centering
\begin{tikzpicture}[node distance=5mm and 6mm,
    every node/.append style={font=\sffamily\footnotesize},
    stage/.append style={minimum width=23mm,minimum height=9mm,inner sep=2pt}]
  % front end, left to right
  \node[stage] (src) {source\\\tiny\texttt{.forge} + prelude};
  \node[stage,right=of src] (lex) {Lexer\\\tiny\texttt{lex.ml}};
  \node[stage,right=of lex] (parse) {Parser\\\tiny\texttt{parser.ml}};
  \node[stage,right=of parse] (expand) {Expansion\\\tiny imports, modules, macros};
  % type system, right to left
  \node[stage,below=8mm of expand] (condense) {Condense\\\tiny\texttt{condense.ml}};
  \node[stage,hot,left=of condense] (tc) {Type inference\\\tiny\texttt{typecheck.ml}};
  \node[stage,left=of tc] (exh) {Exhaustiveness\\\tiny warnings};
  \node[stage,hot,below=8mm of tc] (elab) {Elaboration\\\tiny dictionaries};
  % back ends
  \node[stage,cool,below=8mm of exh |- elab.south,anchor=north] (interp) {Interpreter\\\tiny\texttt{ceval.ml}};
  \node[stage,cool,below=8mm of elab] (spec) {Specialise\\\tiny\texttt{specialize.ml}};
  \node[stage,cool,right=of spec] (mono) {Monomorphise\\\tiny\texttt{mono.ml}};
  \node[stage,cool,right=of mono] (lower) {Lower\\\tiny\texttt{lower\_typed.ml}};
  \node[stage,cool,below=8mm of lower] (llvm) {LLVM IR\\\tiny\texttt{llvm\_emit.ml}};
  \node[stage,cool,left=of llvm] (clang) {Clang\\\tiny\texttt{ls\_runtime.c}};
  \node[stage,left=of clang] (exe) {native\\executable};
  \node[stage,below=8mm of interp] (out) {output};
  \draw[flow] (src) -- (lex);
  \draw[flow] (lex) -- (parse);
  \draw[flow] (parse) -- (expand);
  \draw[flow] (expand) -- (condense);
  \draw[flow] (condense) -- (tc);
  \draw[flow] (tc) -- (exh);
  \draw[flow] (tc) -- (elab);
  \draw[flow] (elab.south) -- ++(0,-4mm) -| (interp.north);
  \draw[flow] (elab) -- (spec);
  \draw[flow] (spec) -- (mono);
  \draw[flow] (mono) -- (lower);
  \draw[flow] (lower) -- (llvm);
  \draw[flow] (llvm) -- (clang);
  \draw[flow] (clang) -- (exe);
  \draw[flow] (interp) -- (out);
\end{tikzpicture}
\caption{The pipeline: tokens, a syntax tree, the core language, then
either the interpreter or native code. Orange: the type system; blue: the
two back ends.}
\label{fig:pipeline}
\end{figure}

\section{Front end}

\paragraph{Lexing and parsing} The lexer turns the source into tokens, each
with its position. The parser is written with parser combinators over the
token list; it backtracks, and on failure reports the token furthest into
the file that it could not get past, with its line. The prelude is parsed
separately so that errors in a program point into the program.

\paragraph{Expansion} \fc|import| inserts the imported files (each once,
rejecting cycles), modules are flattened into qualified names, and macro
invocations are replaced by their expansions.

\paragraph{Condensation} The surface syntax is reduced to a small core
language: operators become applications, blocks become nested \fc|let|s,
functions with several parameters become curried lambdas, a comprehension
guard becomes a generator over a one-element list, and each \fc|impl|
becomes a \emph{dictionary} --- a record of its methods.

\section{Type inference}
\label{sec:impl-inference}

Inference is the Damas--Milner algorithm~\cite{damas-milner}, in two phases. \emph{Constraint generation} walks
an expression and returns a type --- often a fresh type variable --- with
equations between types; \emph{solving} unifies the equations into a
substitution. At each \fc|let|, the equations so far are solved and the type
of the bound names is generalised over the variables that do not occur in
the environment. Figure~\ref{fig:inference} follows the inference of
\fc|let twice f x = f (f x)|.

\begin{figure}[ht]
\centering
\begin{tikzpicture}[font=\small,
    eq/.style={draw=ForgeSteel!60,fill=ForgeSteel!6,rounded corners=2pt,inner sep=4pt,align=left},
    res/.style={draw=ForgeEmber,fill=ForgeEmber!8,rounded corners=2pt,inner sep=5pt,align=center}]
  \node[stage,minimum width=60mm] (e) {\fc|fn f -> fn x -> f (f x)|};
  \node[eq,below left=8mm and -8mm of e] (gen) {%
    $f : \alpha \qquad x : \beta$\\
    $f\ x : \gamma$ \quad gives \quad $\alpha = \beta \to \gamma$\\
    $f\ (f\ x) : \delta$ \quad gives \quad $\alpha = \gamma \to \delta$\\
    result: $\alpha \to \beta \to \delta$};
  \node[eq,right=14mm of gen] (solve) {%
    unify $\beta \to \gamma$ with $\gamma \to \delta$:\\
    $\beta = \gamma$, \quad $\gamma = \delta$\\[2pt]
    substitution: $\alpha \mapsto \delta \to \delta$,\\
    $\beta \mapsto \delta$, \quad $\gamma \mapsto \delta$};
  \node[res,below=8mm of $(gen.south)!0.5!(solve.south)$] (ty) {%
    $(\delta \to \delta) \to \delta \to \delta$ \quad generalised:\quad
    $\forall a.\ (a \to a) \to a \to a$};
  \draw[flow] (e.south) -- node[data,above left]{generate} (gen.north);
  \draw[flow] (gen.east) -- node[data,above]{solve} (solve.west);
  \draw[flow] (solve.south) -- (ty.north east);
  \draw[flow] (gen.south) -- (ty.north west);
\end{tikzpicture}
\caption{Inferring the type of \fc|twice|: generate equations, solve them, generalise.}
\label{fig:inference}
\end{figure}

The solver handles the extensions of the type language: a type-constructor
variable applied to arguments (\fc|f<a>| in \fc|Functor|) unifies with
\fc|Option<Int>| or \fc|List<Int>| by binding the variable to the
constructor; records unify field by field, and an open record's \emph{row}
variable~\cite{remy} absorbs the fields the other side has in addition. Trait predicates
(\fc|Show a|) are collected along the way; at generalisation, those on the
definition's own type variables become part of its scheme, and those on
concrete types must have an instance.

\section{Traits}
\label{sec:impl-traits}

Traits are implemented by \emph{dictionary passing}~\cite{wadler-blott}.
After a definition has been typed, it is typed again with logging switched
on, recording the predicates instantiated at every use of a method or of a
constrained function. Under the definition's solution, each predicate is
resolved: on one of the definition's own type variables, to a dictionary
parameter; on a concrete type, to an instance dictionary, applied to the
dictionaries its \fc|requires| clause needs. The result is an ordinary
program in which methods are record fields and constrained functions take
dictionaries as extra arguments.

The interpreter runs that program as it is. The native back end goes one
step further (\texttt{specialize.ml}): since every dictionary in a whole
program is ultimately a known instance, it makes one copy of each
constrained function per tuple of dictionaries it is used with, and replaces
each method of a known dictionary by that instance's implementation. No
dictionaries remain at run time.

\section{Exhaustiveness}
\label{sec:impl-exhaustiveness}

The checker in \texttt{exhaustiveness.ml} implements Maranget's
\emph{usefulness} algorithm~\cite{maranget}. A vector of patterns is useful
with respect to a matrix of pattern rows if some value matches the vector
and no row. A \fc|case| with arms $p_1, \dots, p_n$ is exhaustive when the
wildcard is not useful with respect to all of them, and arm $p_i$ is
redundant when it is not useful with respect to $p_1, \dots, p_{i-1}$. The
algorithm works column by column, using the scrutinee's type to know the
constructors of each column; when it finds a useful wildcard it also
returns a value that no arm matches, which the warning shows.

\section{The interpreter}

\texttt{ceval.ml} evaluates the core language directly, following the rules
of Section~\ref{sec:dynamic-semantics}: environments map names to values,
functions are closures that capture their environment, and recursive
functions refer to themselves through a mutable reference in their closure.

\section{The native back end}
\label{sec:native}

\paragraph{Monomorphisation} \texttt{mono.ml} type-checks the specialised
program once more, now recording the type of every expression. Under the
solution of each definition, and the instantiation of the copy being built,
every node has a concrete type. Polymorphic functions and values --- at top
level or local --- are copied once per type they are used at, starting from
the program's top-level effects; every local variable gets a unique name. The
result is a typed, monomorphic tree (\texttt{tast.ml}).

\paragraph{Lowering} \texttt{lower\_typed.ml} translates that tree into Min
IR, a small SSA intermediate representation with basic blocks and $\phi$
nodes:
\begin{itemize}
  \item a named function --- top-level or local --- is \emph{lambda-lifted}
    into a function that takes its captured variables as extra parameters,
    and a call that supplies all its arguments is a direct call;
  \item every function value is a curried closure: a code pointer and an
    environment, applied to one argument at a time;
  \item \fc|case| compiles to a chain of tests that jump to the next arm on
    failure, and to a run-time error after the last;
  \item ranges and comprehensions are loops; structural comparison uses one
    generated function per type.
\end{itemize}

\paragraph{Code generation} \texttt{llvm\_emit.ml} prints Min IR as LLVM IR,
and Clang compiles it and links it with the C runtime
(\texttt{runtime/ls\_runtime.c}), which provides strings, the variant and
closure layouts and the run-time errors. Memory is managed by the
Boehm--Demers--Weiser conservative garbage collector, which finds pointers on
the stack, in registers and in global variables without help from the
compiler; strings are allocated as pointer-free (``atomic'') blocks.
Table~\ref{tab:repr} shows how values are represented.

\begin{table}[ht]
\centering\small
\caption{Representation of values in native code.}
\label{tab:repr}
\begin{tabular}{@{}ll@{}}
\toprule
\textbf{Forge type} & \textbf{Native representation} \\
\midrule
\tc{Int}, \tc{Char} & 64-bit integer \\
\tc{Bool} & 1-bit integer \\
\tc{Float} & IEEE double \\
\tc{String} & pointer to a NUL-terminated string \\
\tc{List<a>} & null, or a pointer to a cell \{head, tail\} \\
tuples, records & LLVM structures (records in field-name order) \\
sum types & pointer to \{tag, pointer to the boxed payload\} \\
functions & pointer to \{code, environment\} \\
\bottomrule
\end{tabular}
\end{table}

\section{Testing}

The test suites under \texttt{test/} exercise every stage, from the lexer's
handling of escapes to the error messages of the command-line tools. Two
kinds of tests guard the agreement of the back ends: \emph{differential}
tests run every example program --- including those of this manual --- in
both and require identical output, and a \emph{fuzzer} generates random
well-typed programs (deterministically, from a seed) and does the same.
